\documentclass[11pt]{article}
\usepackage[margin=2cm]{geometry}
\usepackage{amsmath,amssymb,amsthm}
\usepackage{microtype}
\usepackage{enumitem}
\usepackage{booktabs}
\usepackage{xcolor}
\usepackage{tcolorbox}
\tcbuselibrary{skins,breakable}
\usepackage{hyperref}
\hypersetup{colorlinks=true,linkcolor=blue!50!black,urlcolor=blue!50!black}

\newtcolorbox{slide}[1]{
  colback=blue!3,colframe=blue!40!black,
  title={#1},fonttitle=\bfseries,
  breakable, enhanced, before skip=12pt, after skip=12pt
}
\newtcolorbox{paper}{
  colback=blue!5,colframe=blue!50!black,
  title={\textbf{What the article says}},fonttitle=\bfseries,
  breakable, enhanced
}
\newtcolorbox{apply}{
  colback=green!4,colframe=green!40!black,
  title={\textbf{How this applies to our PINN}},fonttitle=\bfseries,
  breakable, enhanced
}
\newtcolorbox{say}{
  colback=yellow!6,colframe=orange!50!black,
  title={\textbf{What to say (rough script)}},fonttitle=\bfseries,
  breakable, enhanced
}
\newtcolorbox{depth}{
  colback=gray!7,colframe=gray!50!black,
  title={\textbf{Depth (for your own understanding)}},fonttitle=\bfseries,
  breakable, enhanced
}
\newtcolorbox{anticipate}{
  colback=red!3,colframe=red!50!black,
  title={\textbf{Likely examiner questions}},fonttitle=\bfseries,
  breakable, enhanced
}

\title{\textbf{Oral exam prep --- Special curriculum}\\
\large Presenting the two articles in depth, with our PINN work as the worked example}
\author{Aleksander Sekkelsten}
\date{}

\begin{document}
\maketitle

\section*{Approach (30-minute talk)}

The talk is a \emph{tour of the two articles}, in their own logical order, with our 2D
quantum-dot PINN as the running worked example. Two papers, two halves, then a synthesis.

\begin{itemize}[leftmargin=*]
\item \textbf{Paper A --- Berner, Grohs, Kutyniok, Petersen, \emph{The Modern Mathematics of
Deep Learning}.} Surveys the field around four questions: (Q1) why does over-parameterization
still generalize, (Q2) what does depth buy in expressivity, (Q3) why do networks beat the curse
of dimensionality, (Q4) why does non-convex optimization work --- plus a section on architectural
priors and one on the natural sciences (PDEs, Schr\"odinger).
\item \textbf{Paper B --- De~Ryck \& Mishra, \emph{Numerical Analysis of PINNs}.} Takes one method
and gives it a numerical-analysis treatment in four matching questions: approximation error,
stability, generalization (quadrature), and \textbf{training}. Its punchline: training difficulty
is governed by the condition number of the \emph{Hermitian square} $L^\ast L$ of the differential
operator, and almost every successful PINN trick is really \emph{preconditioning}.
\end{itemize}

\textbf{The single thread tying the whole talk together:} \emph{the many-body problem is
nominally infinite-dimensional (FCI/HF expand into an exponentially large Hilbert space), but the
physical ground state lives on a low-dimensional manifold, and a neural network is a machine for
finding that manifold --- provided the problem is conditioned well enough for the optimizer to
reach it.} Approximation theory says the manifold exists (Barron, depth, PDE-regularity);
the conditioning theory says whether we can train to it.

\textbf{Slide budget: 18 main + backups. About 90--100\,s per slide.}

\vspace{1em}\hrule\vspace{1em}

% ============================================================
\begin{slide}{Slide 1 --- Outline (45\,s)}
\textbf{On the slide:} two papers; two halves; a one-line statement of our system
(2D parabolic quantum dots, $N\in\{2,6,12,20\}$, $\omega\in\{1,\dots,10^{-3}\}$,
Slater--Jastrow--backflow PINN); and the closing thread (manifold + conditioning).
\end{slide}

\begin{say}
``I'll present the two papers in two halves. First the modern-mathematics review --- approximation,
optimization, generalization, and architecture. Then the De~Ryck--Mishra numerical analysis of
PINNs, which is more focused and whose central result is about conditioning. Throughout I'll use
our 2D quantum-dot calculations as the worked example, and I'll close on the one idea that connects
both papers to our results: the many-body wavefunction lives on a low-dimensional manifold, neural
networks are good at finding such manifolds, and whether we can actually train to it is a question
of how well the problem is conditioned.''
\end{say}

% ============================================================
\begin{slide}{Slide 2 --- Why machine learning for the many-body problem at all (90\,s)}
\textbf{On the slide:}
\begin{itemize}[leftmargin=*]
\item The many-body Hilbert space is exponentially large: FCI expands $\Psi$ in a basis whose
size grows combinatorially with $N$ and the number of orbitals; HF collapses it to a single
determinant and discards correlation.
\item But the physical ground state is \emph{not} a generic vector in that space --- it is smooth,
symmetric, short-range-correlated, and concentrated near a low-dimensional set of configurations.
\item A neural network does not expand in a fixed basis; it \emph{learns} a parameterization that
adapts to the intrinsic structure. This is the core reason ML is attractive here.
\item Empirically (our work): the learned correlator's feature space has \emph{effective rank
$r_{\rm eff}\lesssim 3$} across the whole $(N,\omega)$ grid, even at $N{=}20$ (40-dimensional
configuration space). The network finds a 1--3 dimensional order-parameter manifold.
\end{itemize}
\end{slide}

\begin{say}
``Before either paper, the motivating question: why use a neural network for a quantum many-body
ground state at all? The traditional methods --- full configuration interaction, Hartree--Fock ---
work by expanding the wavefunction in a basis. FCI's basis grows combinatorially; it is exact in
principle and intractable in practice beyond a handful of particles. Hartree--Fock truncates to a
single determinant and throws away correlation. Both are fighting the dimensionality of the Hilbert
space head-on.

The key realization, which both papers articulate in different languages, is that the physical
ground state does not use that whole space. It is smooth, it respects permutation and rotational
symmetry, its correlations are short-ranged, and it concentrates on a low-dimensional set of
configurations. A neural network is precisely a tool for discovering and parameterizing that
low-dimensional structure rather than expanding against it. In our own results we can measure this
directly: the correlator's internal feature space has an effective rank below three across the
entire phase diagram, even for twenty electrons in a forty-dimensional configuration space. The
network compresses an apparently astronomical problem onto a one-to-three dimensional manifold.
That compression is the whole game, and it is what the modern-mathematics article calls escaping
the curse of dimensionality.''
\end{say}

\begin{depth}
\textbf{This frame pays off three times.} (1) In Paper A's approximation theory: Barron spaces and
PDE-regularity are exactly the structural assumptions under which the curse is provably broken ---
they are formal statements of ``the manifold exists.'' (2) In Paper A's architecture section:
equivariance is how we tell the network where the manifold's symmetries are, so it does not waste
capacity. (3) In Paper B's training analysis: even if the manifold exists, gradient descent only
reaches it quickly if the problem is well-conditioned. So ``manifold exists'' (approximation) and
``manifold reachable'' (conditioning) are the two halves, and our thesis is largely about the
second.
\end{depth}

\vspace{2em}\hrule
\begin{center}\Large\textbf{Part I --- The Modern Mathematics of Deep Learning}\end{center}
\hrule\vspace{1em}

% ============================================================
\begin{slide}{Slide 3 --- Paper A: the four questions and the bias--variance puzzle (75\,s)}
\textbf{On the slide:}
\begin{itemize}[leftmargin=*]
\item Classical learning theory: risk $=$ approximation error $+$ estimation (generalization) error;
controlling generalization requires limiting hypothesis-class complexity (VC, Rademacher).
\item Deep networks violate the classical wisdom: they have far more parameters than data, interpolate
the training set, yet generalize. Classical bounds are vacuous.
\item The article's four questions: (Q1) generalization of over-parameterized nets; (Q2) role of
depth; (Q3) beating the curse of dimensionality; (Q4) why optimization succeeds.
\item Section map: 1.2 generalization, 1.3 depth/expressivity, 1.4 curse of dimensionality,
1.5 optimization, 1.6 architectures, 1.7 learned features, 1.8 natural sciences (PDEs/Schr\"odinger).
\end{itemize}
\end{slide}

\begin{say}
``The modern-mathematics article opens by setting up classical statistical learning theory --- the
decomposition of risk into approximation error and estimation error --- and then immediately points
out that deep learning breaks it. Classical theory says: to generalize, restrict the complexity of
your hypothesis class so that fitting the training data means something. Deep networks do the
opposite. They have orders of magnitude more parameters than training points, they interpolate the
data exactly, and they still generalize to new data. Every classical bound --- VC dimension,
Rademacher complexity --- is vacuous in this regime. So the article organizes the modern theory
around four questions: why over-parameterized networks generalize, what depth buys you, how networks
beat the curse of dimensionality, and why a non-convex optimization problem is solvable at all. I'll
take these roughly in the article's order, and tie each to our work.''
\end{say}

% ============================================================
\begin{slide}{Slide 4 --- Approximation I: universal approximation is not enough; the curse (90\,s)}
\textbf{On the slide:}
\begin{itemize}[leftmargin=*]
\item Universal approximation (Cybenko 1989, Hornik 1991): a one-hidden-layer network with
non-polynomial activation is dense in $C(K)$. A \emph{topological} statement --- no rate.
\item With rates, the curse appears: approximating a $C^s$ function on $\mathbb R^d$ to accuracy
$\varepsilon$ needs $\sim\varepsilon^{-d/s}$ parameters. Exponential in $d$.
\item \textbf{Dimensional analysis of our problem.} $d_{\rm config}=Nd$: $N{=}2\to4$,
$N{=}6\to12$, $N{=}12\to24$, $N{=}20\to40$. A naive $C^s$ rate at $d{=}40$ is hopeless
($\varepsilon^{-40/s}$).
\item Escape routes the article names: smoothness, compositionality, low-dimensional support
(manifold), \textbf{symmetry}, and \textbf{Barron regularity} (next slide).
\end{itemize}
\end{slide}

\begin{say}
``The first pillar is approximation. The theorem everyone remembers is universal approximation: one
hidden layer suffices to approximate any continuous function on a compact set. The article is
careful to call this what it is --- a topological density statement with no rate attached. The
moment you ask how many neurons you need for accuracy $\varepsilon$, the curse of dimensionality
appears: for a function with $s$ derivatives on a $d$-dimensional domain, you need on the order of
$\varepsilon^{-d/s}$ parameters.

For our problem this is the central obstacle. The configuration space is $Nd$-dimensional: four
dimensions for two electrons, but forty for twenty electrons in 2D. A naive smoothness-based rate at
$d{=}40$ is utterly hopeless --- $\varepsilon^{-40/s}$ is astronomical for any realistic $s$. So the
entire question of whether a neural wavefunction is feasible reduces to: does our target have extra
structure that breaks the curse? The article lists the escape routes --- smoothness alone is not
enough, but compositional structure, low-dimensional support, symmetry, and what's called Barron
regularity all are. Our ansatz is engineered to exploit every one of these.''
\end{say}

\begin{depth}
\textbf{Why FCI/HF do not escape and NN can.} FCI is essentially a linear approximation in a fixed
(if huge) basis --- it is a spectral method, and spectral/finite-element/sparse-grid methods
provably suffer the curse for generic high-dimensional targets (the article states this explicitly
in 1.4.3). HF is a rank-one nonlinear approximation that discards correlation. Neural networks are
\emph{nonlinear, adaptive} approximators: the basis functions move during training to where the
target has mass. The article's depth and Barron results are the rigorous statements that this
adaptivity can beat the curse.
\end{depth}

% ============================================================
\begin{slide}{Slide 5 --- Approximation II: Barron spaces and dimension-independent rates (100\,s)}
\textbf{On the slide:}
\begin{itemize}[leftmargin=*]
\item \textbf{Barron's theorem (1.27).} If $g$ has a finite first Fourier moment,
$C_g=\int_{\mathbb R^d}\|\xi\|_2\,|\hat g(\xi)|\,d\xi<\infty$, then a two-layer network with $n$
neurons approximates $g$ in $L^2$ with error $\le c\,C_g/\sqrt n$.
\item The rate $n^{-1/2}$ is \textbf{independent of dimension $d$}. The price of dimension is hidden
in $C_g$, not in the exponent.
\item Mechanism: write $g$ via its inverse Fourier transform as an expectation, then approximate the
expectation by Monte Carlo over $n$ random neurons (law of large numbers). The $n^{-1/2}$ is the MC
rate.
\item \textbf{Generalized / spectral Barron spaces} (E, Wojtowytsch; Barron--Klusowski) extend this
to deep, compositional functions --- and crucially to PDE solutions, including \emph{static
Schr\"odinger equations in spectral Barron spaces} (Chen--Lu--Lu--Zhou 2023, cited by Paper B).
\end{itemize}
\end{slide}

\begin{say}
``This is the result the article uses to break the curse, and it is the one I'd want to be able to
state precisely. Andrew Barron showed in 1992--93 that if a function's Fourier transform has a finite
first moment --- the integral of the frequency magnitude times the spectrum is finite --- then a
two-layer network with $n$ neurons approximates it with $L^2$ error proportional to $1/\sqrt n$. The
remarkable thing is that the exponent does not depend on the dimension at all. The dimension
dependence is pushed into the constant $C_g$, the Barron norm. So for functions in this class, the
curse of dimensionality in the \emph{rate} is gone.

The proof idea is beautiful and physical: you write the function via its inverse Fourier transform as
an integral, reinterpret that integral as an expectation over a probability distribution on
frequencies, and then approximate the expectation by drawing $n$ random neurons. The $1/\sqrt n$ is
just the Monte Carlo rate from the law of large numbers --- each neuron is one sample. So a wide
network is doing Monte Carlo integration in frequency space.

Why this matters for us: the theory was later generalized --- E and Wojtowytsch, Barron and
Klusowski --- to deep and compositional Barron spaces, and then to PDE solutions. There is now a
regularity theory for the \emph{static Schr\"odinger equation in spectral Barron spaces}, which
Paper B cites. So the claim that our wavefunction is efficiently approximable by a neural network is
not just hope --- it sits inside an existing approximation-theoretic framework.''
\end{say}

\begin{depth}
\textbf{The connection to conditioning (preview of Part II).} The Barron norm weights the Fourier
tail by $\|\xi\|$. Apply a differential operator and you multiply the spectrum by powers of
$\|\xi\|$, so the residual of a PDE lives in a \emph{worse} Barron space than the solution: the
derivative amplifies high frequencies. This is the single most important bridge between the two
papers --- it is why a strong-form residual loss is harder to approximate \emph{and} harder to
condition than the solution itself, and why our analytic cusp (which removes the high-frequency
content the network would otherwise have to represent) is doing approximation-theoretic work, not
just numerical hygiene. We return to this on slides 13--15.

\textbf{The multi-electron Schr\"odinger equation appears in Paper A itself.} Section 1.4.3 cites
Al-Hamdani et al.\ (2020): two gold-standard methods disagree on interaction energies for large
delocalized molecules because classical numerical representations are not expressive enough for
high-dimensional wavefunctions with long-range interactions. The article then cites results showing
NNs provably avoid the curse for elliptic PDEs, Hamilton--Jacobi--Bellman, parametric PDEs, etc.
Our work is an empirical instance of exactly this claim.
\end{depth}

% ============================================================
\begin{slide}{Slide 6 --- Approximation III: depth, compositionality, and non-smooth obstructions (75\,s)}
\textbf{On the slide:}
\begin{itemize}[leftmargin=*]
\item Depth separation (Telgarsky; Eldan--Shamir): functions a depth-$L$ network represents with
$O(L)$ neurons need $\Omega(2^L)$ neurons at constant depth. Depth is qualitatively, not
incrementally, more expressive.
\item Compositionality (1.4.1): NNs efficiently represent compositions of low-dimensional functions;
local-coordinate / manifold structure reduces effective dimension $d\to d'\ll d$.
\item \textbf{Non-smooth obstruction in our problem:} the Kato cusp. $\partial_r\log\Psi$ has a
\emph{kink} at coalescence; the wavefunction is not $C^1$ there. A smooth (or Barron-regular)
network needs unbounded spectral budget to represent it $\Rightarrow$ we hard-wire it analytically.
\end{itemize}
\end{slide}

\begin{say}
``Two more approximation results, then the obstruction. First, depth separation: there are functions
a deep network represents with linearly many neurons that a shallow network needs exponentially many
to match. Depth is not a tuning knob; it gives access to a categorically larger function class at
fixed budget. Second, compositionality: the article shows networks efficiently represent compositions
of low-dimensional maps, which is how they exploit manifold structure --- they reduce the effective
dimension from the ambient $d$ to the intrinsic $d'$.

But there is an obstruction specific to our physics that no amount of depth fixes: the cusp. The
Coulomb interaction forces the exact wavefunction to have a kink --- a discontinuous radial
derivative --- wherever two electrons meet. That is a non-smooth feature; in Barron or Sobolev terms
it has slowly decaying high-frequency content. A smooth network can only approximate a kink with
ever-growing weights. So rather than spend the network's capacity fighting a singularity it can never
exactly represent, we hard-wire the cusp with an analytic term and let the network learn only the
smooth remainder. This is the approximation-theory justification for a choice that looks like mere
engineering.''
\end{say}

% ============================================================
\begin{slide}{Slide 7 --- Optimization: loss landscape, NTK, spectral bias, implicit bias (100\,s)}
\textbf{On the slide:}
\begin{itemize}[leftmargin=*]
\item The loss is non-convex; no global guarantee exists. The article's resolution is the
\textbf{wide-network / kernel regime}.
\item \textbf{Neural Tangent Kernel:} $\Theta(x,x')=\langle\nabla_\theta f_\theta(x),
\nabla_\theta f_\theta(x')\rangle$. As width $\to\infty$, $\Theta$ is constant along training and the
dynamics linearize --- training becomes kernel regression.
\item \textbf{Spectral bias:} components along large-eigenvalue NTK eigenvectors are learned fast,
small-eigenvalue (high-frequency) components slow.
\item \textbf{Implicit bias:} among many interpolating minima, GD selects the minimum-norm one ---
generalization without explicit regularization.
\item Our optimizer: \textbf{Stochastic Reconfiguration} $=$ natural gradient with the quantum Fisher
matrix $S_{ij}=\mathrm{Cov}[O_i,O_j]$, $O_i=\partial_{\theta_i}\log\Psi$. $S$ is the QMC NTK.
\end{itemize}
\end{slide}

\begin{say}
``The optimization pillar. The loss surface of a deep network is non-convex with exponentially many
critical points, and there is no worst-case guarantee that gradient descent finds a good minimum.
The article's resolution is the wide-network limit and the Neural Tangent Kernel. The NTK is the
inner product of parameter-gradients of the network output at two inputs. Jacot's result is that as
the width goes to infinity, this kernel stays constant during training, the dynamics linearize, and
training reduces to kernel regression with the NTK --- a convex problem whose convergence is governed
by the kernel's eigenvalue spectrum.

Two phenomena drop out. Spectral bias: directions with large NTK eigenvalues are learned quickly,
small-eigenvalue directions slowly, and for smooth kernels the small eigenvalues correspond to high
frequencies --- so networks fit low frequencies first. Implicit bias: of the infinitely many
parameter settings that interpolate the data, gradient descent prefers the minimum-norm one, which
is why over-parameterized nets generalize.

Our optimizer makes this concrete. We use Stochastic Reconfiguration, which is natural gradient with
the quantum Fisher information matrix as preconditioner. The Fisher matrix is built from the
log-derivatives of the wavefunction --- it is, up to the sampling measure, exactly the NTK of the
log-amplitude. So when SR accelerates convergence, the reason is that it rescales the slow,
small-eigenvalue NTK directions that spectral bias would otherwise leave untrained. This is the same
mechanism Paper B will formalize as preconditioning.''
\end{say}

\begin{depth}
\textbf{Where the wide-network story breaks for us, and why that is the interesting part.} The NTK
constancy requires very wide networks; ours are deliberately compact ($\sim 50$k correlator
parameters). So we are \emph{not} cleanly in the kernel regime, and the feature manifold genuinely
moves during training --- which is exactly why the correlator can \emph{reorganize} which branch
dominates as $\omega$ changes (slide 17). The kernel-regime theory is the idealization; our regime is
the feature-learning regime where the manifold itself adapts. Both papers acknowledge this gap: the
kernel regime is where the theory is clean, the feature-learning regime is where the empirical wins
are. We live in the second.
\end{depth}

% ============================================================
\begin{slide}{Slide 8 --- Generalization: kernel, norm bounds, double descent --- and our analogue (90\,s)}
\textbf{On the slide:}
\begin{itemize}[leftmargin=*]
\item Article's candidates (1.2): kernel-regime generalization, norm/margin-based bounds (control the
\emph{size} of weights, not their count), implicit regularization, and \textbf{double descent}
(test error falls again past the interpolation threshold).
\item Honest verdict: phenomenology is robust, a tight theory is not yet settled.
\item \textbf{Our analogue:} no train/test split. The generalization measure is the
\textbf{variance of the local energy}, $\mathrm{Var}[E_L]$. Zero-variance principle:
$\mathrm{Var}[E_L]=0\iff\Psi_\theta$ is an exact eigenstate.
\item Sample efficiency: we reach DMC-level energies with $\sim 3\times10^3$ samples per SR step,
$\sim 10^3$ steps --- orders of magnitude fewer than typical neural-VMC.
\end{itemize}
\end{slide}

\begin{say}
``Generalization is the pillar where the article is most candid that the theory is incomplete. It
surveys the serious candidates: generalization in the kernel regime; norm-based and margin-based
bounds, which control the magnitude of the weights rather than their number and so survive
over-parameterization; implicit regularization again; and double descent, the striking empirical
curve where increasing capacity past the interpolation threshold makes the test error fall a second
time, contradicting the classical bias--variance U-curve. The honest summary is that we have robust
phenomenology and no fully settled theory.

For our problem the classical notion does not even apply --- there is no held-out test set. The right
generalization measure is the variance of the local energy over configurations. The zero-variance
principle says this variance is zero exactly when the wavefunction is an exact eigenstate, so it is
simultaneously a measure of accuracy and of generalization. Our practical signature is sample
efficiency: we reach diffusion-Monte-Carlo-level energies with a few thousand samples per step and
about a thousand steps, orders of magnitude fewer samples than comparable neural approaches. That
efficiency is the implicit-bias story in physical form: a symmetry-restricted, well-conditioned
ansatz needs far less data to pin down the correlator.''
\end{say}

% ============================================================
\begin{slide}{Slide 9 --- Architectures: from CNN equivariance, via ResNet--ODE, to GNN and CTNN (100\,s)}
\textbf{On the slide:}
\begin{itemize}[leftmargin=*]
\item \textbf{CNN (1.6.1):} three principles --- local receptive fields, parameter sharing,
equivariant representations. Restricting to convolutions $=$ restricting the hypothesis class to
translation-equivariant maps. Universality holds on that restricted class.
\item \textbf{ResNet (1.6.2):} skip connections $\bar\Phi^{(\ell)}=\varrho(\Phi^{(\ell)})+
\bar\Phi^{(\ell-1)}$ are an Euler discretization of an ODE $\dot\varphi=h(t,\varphi)$ ---
``continuous-time'' networks; cures vanishing/exploding gradients.
\item \textbf{GNN (explicitly flagged as omitted, 1.6):} non-Euclidean inputs; permutation
equivariance; cited to Bronstein, Wu.
\item \textbf{Generalization chain:} CNN (translation) $\to$ GNN (permutation, graph) $\to$
\textbf{CTNN} $=$ a continuous-time / ResNet-style message-passing GNN. This is \emph{our} backflow.
\end{itemize}
\end{slide}

\begin{say}
``Now the architecture section, which is where the article stops being generic learning theory and
starts being directly about our network. The article's strongest positive claim is that universal
approximation alone does not beat the curse --- \emph{structure} does --- and architecture is how you
encode structure. It develops this for convolutional networks: three principles --- local receptive
fields, parameter sharing, and equivariant representations. Restricting a network to convolutions is
exactly restricting its hypothesis class to translation-equivariant functions, and you can prove
universality \emph{on that restricted class}. You get expressivity and reduced effective dimension
simultaneously.

Then ResNets: the skip connection turns out to be an Euler step of an ordinary differential equation,
so a deep residual network is a discretized continuous-time flow. This is both how they cure the
vanishing-gradient problem and how you get the ``neural ODE'' viewpoint.

The article explicitly flags that it omits graph neural networks --- the generalization of CNNs to
non-Euclidean, permutation-symmetric inputs --- and points to the Bronstein and Wu surveys. That
omission is exactly the gap our architecture fills. The chain is: CNN encodes translation
equivariance; GNN generalizes to permutation equivariance on a graph of interacting objects; and our
CTNN --- continuous-time neural network --- is a ResNet-style, ODE-limit message-passing GNN. So our
backflow is the natural composition of two ideas the article develops separately: graph
message-passing and the continuous-time residual flow.''
\end{say}

\begin{depth}
\textbf{Why ``continuous-time'' is the right name and the right design.} The ResNet-as-ODE picture
(1.6.2) says a residual block integrates $\dot\varphi=h(t,\varphi)$. Our CTNN backflow produces a
displacement field $\Delta_\beta$ as a bounded, near-identity flow: zero-initialized last layer,
$\tanh$ output, learnable positive scale, so $\Delta_\beta\approx 0$ at initialization and grows
smoothly --- exactly an ODE flow started at the identity. This is why its displacements stay small
($\|\Delta x\|/\ell\lesssim 0.1$, never blowing up like the conventional BackflowNet's $0.95$): the
continuous-time parameterization is an \emph{implicit regularizer}, the architectural analogue of the
implicit bias from slide 7. The article's two separate observations (ResNets are ODEs; architecture
encodes priors) combine into one design decision in our network.
\end{depth}

% ============================================================
\begin{slide}{Slide 10 --- Our architecture in detail: three networks $+$ backflow, and why (100\,s)}
\textbf{On the slide:}
\[
\Psi_{\theta,\beta}(R)=\mathrm{SD}\big(\tilde R+\Delta_\beta(\tilde R)\big)\,
\exp\Big(\underbrace{\textstyle\sum_{i<j}u_{\sigma_i\sigma_j}(\tilde r_{ij})}_{\text{analytic cusp}}
+\underbrace{W_\theta(\tilde R)}_{\text{neural correlator}}\Big)
\]
\begin{itemize}[leftmargin=*]
\item \textbf{Three correlator networks.} $\phi$: particle branch (Deep Sets), per-particle embedding
$\to$ symmetric pool. $\psi$: pair branch, embeds every pair $(i,j)$ from \emph{safe} features $\to$
pool. $\rho$: readout, maps $[\bar h^\phi\,\|\,\bar h^\psi\,\|\,\mathbf g]\to$ scalar.
\item \textbf{Why Deep Sets alone fails:} particle-only pooling discards relative geometry --- cannot
represent two-body correlation. Pair branch is mandatory (Haas: Deep Sets needs correlation terms).
\item \textbf{Why so many inputs?} 6 soft-core radial features $+$ raw $(\Delta x,\Delta y)$ $+$
one-body $(x_i,y_i)$ $+$ spin $+$ 2 global scalars. Redundancy gives the optimizer well-conditioned,
bounded channels to choose from.
\item \textbf{Backflow networks $\phi_\beta,\psi_\beta$:} the CTNN message-passing flow on the
determinant coordinates only.
\end{itemize}
\end{slide}

\begin{say}
``Here is our actual ansatz and the design logic. The wavefunction is a Slater determinant times the
exponential of an analytic cusp plus a neural correlator, with the determinant optionally evaluated
on backflow-displaced coordinates. The neural correlator is built from three hand-crafted networks.
The particle branch $\phi$ is a Deep Sets network: it embeds each particle and pools symmetrically,
giving permutation invariance by construction. The pair branch $\psi$ embeds every pair of particles
from carefully designed features and pools. The readout $\rho$ takes the pooled particle embedding,
the pooled pair embedding, and two global scalars, and maps them to the scalar log-amplitude
correction.

Why three networks and not one Deep Sets? Because a particle-only Deep Sets, however deep, literally
cannot represent two-body correlation: once you pool the per-particle embeddings, the relative
geometry is gone. This is a theorem-level limitation, and it is exactly what the Haas result we cite
points out --- Deep Sets needs explicit correlation terms. The pair branch is not an enhancement; it
is the minimum required to capture correlation physics.

Why so many inputs to the pair branch? We feed six soft-core radial features, the raw coordinate
differences, the one-body positions, a spin indicator, and two global statistics. The redundancy is
deliberate and it is a conditioning choice: the soft-core radial features have bounded derivatives at
coalescence, so they are safe channels for the Laplacian, and we give the optimizer a rich, mostly
well-behaved menu to allocate over. The attribution analysis then shows the network does the
sensible thing --- it suppresses the unsafe raw-coordinate channels and leans on the radial features,
and the balance it strikes shifts with the physical regime.''
\end{say}

\begin{depth}
\textbf{Did the extra inputs help, and do they condition better?} Yes, and the attribution data
shows how. At tight confinement ($\omega{=}1$) the radial distance $|r|$ carries $40.5\%$ of the
Jastrow's input attribution; as the trap weakens, \emph{spin} takes over --- $51.4\%$ at $\omega{=}
0.1$, $71.6\%$ at the Wigner limit. The raw $(\Delta x,\Delta y)$ channels, which are \emph{not} safe
at coalescence (their derivatives don't vanish there), carry consistently low attribution: the
network learns to avoid the badly-conditioned inputs on its own. So the many-input design is a
conditioning device --- it offers redundant, bounded ways to express the same correlation, and the
optimizer migrates onto whichever channel is both physically informative \emph{and} numerically safe
in each regime. This is the architectural face of the conditioning thesis from Part II: give the
network well-conditioned coordinates and it finds the manifold faster.

\textbf{Permutation/rotation bookkeeping.} $\phi,\psi$ pool symmetrically $\Rightarrow$ permutation
invariance. Backflow uses vector messages $\mathbf r_{ij}$ and centering $\Rightarrow$ rotation
equivariance and (near) translation invariance. Antisymmetry is the determinant. So every symmetry
the article would have us encode is encoded.
\end{depth}

\vspace{2em}\hrule
\begin{center}\Large\textbf{Part II --- Numerical Analysis of PINNs}\end{center}
\hrule\vspace{1em}

% ============================================================
\begin{slide}{Slide 11 --- Paper B: PINNs as a numerical method; four matching questions (90\,s)}
\textbf{On the slide:}
\begin{itemize}[leftmargin=*]
\item A PINN solves $\mathcal L u=f$ (with BCs) by minimizing the residual in strong / weak /
variational form over a network class. PINN $=$ strong form $+$ neural ansatz.
\item Total error decomposition (Paper B's spine):
\[
\underbrace{\|u_\theta-u\|}_{\text{total}}\ \lesssim\ C_{\rm stab}\big(
\underbrace{\mathcal E_{\rm approx}}_{\text{Q1, \S4}}+
\underbrace{\mathcal E_{\rm gen}}_{\text{Q3 quadrature, \S6}}+
\underbrace{\mathcal E_{\rm opt}}_{\text{Q4 training, \S7}}\big),\quad
C_{\rm stab}\ \text{(Q2, \S5)}.
\]
\item Q1 approximation, Q2 stability, Q3 generalization (quadrature), Q4 training.
\item Verdict (their conclusion, verbatim spirit): \emph{the key bottleneck is training, and it is
governed by the spectral properties of the Hermitian square $L^\ast L$; preconditioning is the cure.}
\end{itemize}
\end{slide}

\begin{say}
``Now the second paper. It is much more focused: one method, full numerical-analysis treatment. A
PINN solves a PDE by minimizing its residual --- in strong, weak, or variational form --- over a
neural ansatz; the vanilla PINN is the strong form with a neural network. De~Ryck and Mishra
decompose the total error, the gap between the trained network and the true solution, into a stability
constant times the sum of three pieces that mirror Paper A's pillars exactly: approximation error,
generalization or quadrature error, and optimization or training error. The stability constant is the
fourth, PDE-specific ingredient.

Their organizing questions are: Q1, is the solution in the network class; Q2, is the PDE stable so
that small residual implies small error; Q3, does the finite-sample loss approximate the true loss;
and Q4, can we actually train. And their headline conclusion --- which I'll quote almost verbatim
later --- is that the binding constraint is Q4, training, and that training difficulty is governed by
the spectral properties of the Hermitian square of the differential operator, with preconditioning as
the remedy. That conclusion is the heart of what I want to connect to our work.''
\end{say}

% ============================================================
\begin{slide}{Slide 12 --- Stability (Q2): small residual $\Rightarrow$ small error, only sometimes (90\,s)}
\textbf{On the slide:}
\begin{itemize}[leftmargin=*]
\item For $\mathcal L u=f$, $C_{\rm stab}=\|\mathcal L^{-1}\|$ on the relevant space. Coercive elliptic
$\Rightarrow$ bounded (Lax--Milgram). Navier--Stokes, conservation laws $\Rightarrow$ delicate; Paper
B proves case-by-case stability.
\item \textbf{Eigenproblem (us):} $\mathcal L=\hat H-E$ is \emph{not invertible} on the eigenspace;
$C_{\rm stab}=1/\text{gap}$ to the next eigenvalue. Strong-form residual is unstable near
degeneracy.
\item \textbf{Fix = change the form.} Rayleigh quotient $E[\Psi]=\langle\Psi|\hat H|\Psi\rangle/
\langle\Psi|\Psi\rangle\ge E_0$. This is the \textbf{Deep Ritz} variant; bounded below, stability
automatic. ``Regularity dictates the form.''
\end{itemize}
\end{slide}

\begin{say}
``Stability is the ingredient that has no analogue in ordinary deep learning. For a generic PDE, a
small residual implies a small solution error only if the operator is invertible with a bounded
inverse --- the stability constant. For coercive elliptic problems this is automatic; for
Navier--Stokes or hyperbolic conservation laws it is subtle, and a large fraction of Paper B is
devoted to proving stability estimates PDE by PDE.

For us the operator is $\hat H - E$, and at the eigenvalue it is by construction singular --- the
stability constant is the inverse of the gap to the first excited state. Near a degeneracy the gap
closes and the strong-form residual stops controlling the error. The fix is the one the article
recommends in general: when the strong form is ill-posed, change the form. We minimize the Rayleigh
quotient, the variational energy, which is bounded below by the true ground-state energy. That is the
Deep Ritz method, and the article treats it as a PINN variant with automatic stability. Their slogan
is `regularity dictates the form,' and our pipeline is a direct instance: strong-form residual only
for cheap pretraining, variational form for the accurate production runs.''
\end{say}

% ============================================================
\begin{slide}{Slide 13 --- Training (Q4): the condition number is the whole story (110\,s) \emph{[Part II centerpiece]}}
\textbf{On the slide:}
\begin{itemize}[leftmargin=*]
\item Linearize GD around $\theta_0$; $q_i=\partial_{\theta_i}u_{\theta_0}$. Convergence governed by
$A_{ij}=\langle\mathcal L q_i,\mathcal L q_j\rangle$.
\item \textbf{Theorem (7.9, 7.21):} steps to accuracy $\varepsilon$ scale as
$\#(\varepsilon)=O\!\big(\kappa(A)\ln\tfrac1\varepsilon\big)$, $\kappa(A)=\lambda_{\max}/\lambda_{\min}$.
\item \textbf{Theorem 7.10:} $\kappa(A)$ equals the condition number of $A\circ TT^\ast$, where
$A=L^\ast L$ is the \textbf{Hermitian square} of the operator and $TT^\ast$ is the \textbf{kernel
integral operator of the NTK}.
\item For the Laplacian $L=-\Delta$: $L^\ast L=\Delta^2$ (bi-Laplacian), Fourier symbol
$|\mathbf k|^4\Rightarrow\kappa\sim(k_{\max}/k_{\min})^4$.
\item \textbf{``Training complexity is entirely due to the spectral properties of $L^\ast L$.''}
Supervised learning is the special case $L=\mathrm{Id}$, $A=$ NTK only.
\end{itemize}
\end{slide}

\begin{say}
``This is the technical core of the second paper and, I think, the single most important slide of the
whole talk. They linearize gradient descent around initialization, define the matrix $A$ whose
entries are inner products of the operator applied to the parameter-gradient features, and prove that
the number of gradient steps to reach accuracy $\varepsilon$ scales linearly with the condition
number of $A$ times a log factor. Slow training is large condition number, full stop.

Then comes the result that names the mechanism. Theorem 7.10 shows the condition number of $A$ equals
the condition number of the Hermitian square of the differential operator --- $L^\ast L$ --- composed
with the kernel integral operator of the neural tangent kernel. So training speed factorizes into two
pieces: the NTK, which is the pure machine-learning part and is all you have in ordinary supervised
learning; and $L^\ast L$, the operator part, which is what makes PINNs special and hard. For the
Laplacian, $L^\ast L$ is the bi-Laplacian, whose Fourier symbol is $|\mathbf k|^4$, so the condition
number scales as the fourth power of the ratio of the largest to smallest resolved frequency. That
quartic is the precise mathematical form of `Laplacian stiffness.'

Their sentence, which I'll basically quote, is that the entire extra difficulty of training a PINN
relative to supervised learning is due to the spectral properties of $L^\ast L$. Supervised learning
is the special case where the operator is the identity and only the NTK remains. Once you see this,
every PINN training trick in the literature is reinterpreted as preconditioning --- as trying to make
$TT^\ast$ approximate the inverse of $L^\ast L$, a Green's function.''
\end{say}

\begin{depth}
\textbf{This is in our appendix, derived independently and citing this paper.} Our Laplacian-
conditioning appendix reproduces exactly this chain: linearize, form $A=J^\ast J$, identify the
operator-level object $L^\ast L\approx\tfrac14\Delta^2$, get the $|\mathbf k|^4$ weight, and conclude
$\kappa\sim(k_{\max}/k_{\min})^4$ with $k_{\max}\sim 1/\ell$ for the smallest length scale $\ell$
(nodal/cusp/localization structure). So the slide is not borrowed --- we have the same theorem in our
own notation, and we use it to explain our regime dependence: small-$\omega$ Wigner structure and
sharp nodes mean small $\ell$, large $k_{\max}$, worse conditioning.

\textbf{The preconditioner-as-Green's-function idea (7.4).} They show preconditioning is a linear
reparameterization $P$ with $\tilde A=P^\top A P$, and the ideal $P$ makes $TT^\ast\approx A^{-1}$.
Their worked example is learnable Fourier features for the Poisson equation: the natural
preconditioner is $P_{kk}=1/k^2$, which exactly undoes the $k^4$ of the bi-Laplacian and collapses
$\kappa$ from $\ge k^4$ to $O(1)$. SR is the QMC version of this $P$.
\end{depth}

% ============================================================
\begin{slide}{Slide 14 --- The conditioning thesis applied: how WE reduced $\kappa$ (110\,s)}
\textbf{On the slide:} every one of our design choices is a preconditioner.
\begin{center}
\begin{tabular}{ll}
\toprule
\textbf{Our technique} & \textbf{What it conditions} \\
\midrule
Trap units $\tilde r=\sqrt\omega\,r$ & rescales $\ell\to O(1)$ across $\omega$; shrinks $k_{\max}/k_{\min}$ \\
Soft-core features $ds/dr\to0$ & removes high-$|\mathbf k|$ content from the learned part \\
Analytic cusp $\gamma r e^{-r}$ & cancels the $1/r$ residual spike (worst high-$\mathbf k$) \\
Variational / REINFORCE form & lowers operator order $\Rightarrow$ avoids $L^\ast L$ stiffness \\
Stochastic Reconfiguration & explicit preconditioner $P$, $TT^\ast\!\approx\!A^{-1}$ (Fisher) \\
\bottomrule
\end{tabular}
\end{center}
\textbf{Negative result we can defend:} K-FAC / Adam preconditioning on collocation \emph{does not}
break the $0.32\%$ error floor, because they condition \emph{parameter space} ($TT^\ast$), not the
\emph{operator} ($L^\ast L$). Only changing the form (VMC$+$SR) reaches $<0.01\%$.
\end{slide}

\begin{say}
``This is the analysis I most want to do, and it follows directly from Theorem 7.10. If training
speed is set by the condition number of $L^\ast L$ composed with the NTK, then every successful
design choice in our thesis should be readable as a preconditioner --- and it is.

Trap units: rescaling coordinates by root-omega fixes the characteristic length to order one across
three decades of confinement, which directly shrinks the ratio of largest to smallest frequency that
the network must resolve --- it attacks $k_{\max}/k_{\min}$ in the quartic. Soft-core features with
vanishing derivative at coalescence remove high-frequency content from the part of the wavefunction
the network has to learn. The analytic cusp cancels the $1/r$ spike in the residual, which is
precisely the worst high-frequency contribution. The variational form lowers the order of the
operator --- this is the `regularity dictates form' lever, trading the stiff strong-form $L^\ast L$
for a milder object. And Stochastic Reconfiguration is the explicit preconditioner: the Fisher matrix
plays the role of $P$, approximating the inverse of $A$, the QMC version of the Fourier-feature
preconditioner in their Poisson example.

The most instructive part is a negative result we can defend rigorously. On the collocation pipeline,
we tried K-FAC and various Adam schedules to fix backflow training, and nothing broke a roughly
0.3-percent error floor --- while the identical ansatz under VMC plus SR reaches below 0.01 percent.
Theorem 7.10 explains why: K-FAC and Adam precondition the parameter-space factor, the NTK piece
$TT^\ast$. But our stiffness is in the operator factor $L^\ast L$ --- it is operator-level, not
parameter-level. You cannot precondition your way out of it in parameter space; you have to change
the form of the problem. That is exactly what switching to the variational principle does, and it is
why the only thing that worked was abandoning the strong-form collocation framework entirely.''
\end{say}

\begin{depth}
\textbf{This negative result is the strongest single point you have.} It is a clean, falsifiable
prediction of Paper B's central theorem, tested in our appendix across $\sim 17$ interventions
(sampling tricks, loss modifications, regularizations, K-FAC), \emph{all} of which fail to beat
$0.32\%$, versus $<0.01\%$ for VMC$+$SR. The mechanism is the self-referential singularity: backflow
parameters move the node, so collocation points that informed the gradient are invalidated by the
step --- whereas VMC resamples from the current $|\Psi_\theta|^2$ every step, co-adapting the measure
with the node. If an examiner asks ``so is preconditioning the answer?'', the precise answer is:
\emph{parameter-space preconditioning is necessary but not sufficient; when the stiffness is in
$L^\ast L$, you must also reduce operator order (change the form). Our data shows both halves.}
\end{depth}

% ============================================================
\begin{slide}{Slide 15 --- Coulomb conditioning: the cusp as a cancellation problem (90\,s)}
\textbf{On the slide:}
\begin{itemize}[leftmargin=*]
\item 2D cusp conditions (spin-resolved): antiparallel $\partial_r\log\Psi|_0=1$; parallel (node)
$\partial_r\log J|_0=\tfrac13$.
\item A cusp-slope error $\delta a$ gives a local-energy spike $E_L\supset -\delta a/r$ --- arbitrarily
small mismatch, arbitrarily large spike.
\item In 2D the squared spike $\int_{r<\varepsilon}(1/r)^2\,d^2r=\infty$ (log-divergent) --- strong
residual is dominated by rare close-pair samples; heavy-tailed MC.
\item $L^\ast L$ for $L=\hat T+(\hat V-E)$ contains $(\hat V-E)^2\sim 1/r^2$ and mixed terms with
$\nabla V$, $\Delta V$ --- even more singular. Coulomb \emph{inflates} the spectral spread.
\item \textbf{Cure (ours):} analytic cusp $+$ soft features make $E_L$ bounded $\Rightarrow$ the
quadrature assumptions of \S6 hold and SR forces become finite-variance.
\end{itemize}
\end{slide}

\begin{say}
``One more conditioning slide, because Coulomb deserves its own. The cusp conditions in 2D are
spin-resolved: for antiparallel spins the log-derivative of the wavefunction at contact must equal
one; for parallel spins the determinant already has a linear node and the Jastrow slope must be
one-third. The danger is sensitivity: if the learned cusp slope is off by a small amount, the local
energy acquires a spike that goes like that error over $r$ --- an arbitrarily small cusp mismatch
produces an arbitrarily large energy spike at coalescence.

In two dimensions this is especially vicious, because the squared spike is log-divergent when you
integrate it against the 2D area element. So a strong-form squared residual is dominated by the rare
samples where two electrons are close, and the Monte Carlo estimator becomes heavy-tailed. In the
operator language of the previous slides, the Hermitian square $L^\ast L$ now contains a $1/r^2$ term
from the potential squared, plus mixed kinetic-potential terms involving gradients and Laplacians of
the Coulomb potential, which are even more singular. Coulomb directly inflates the spectral spread
that controls the condition number.

Our cure is the analytic cusp plus the soft-core features. Together they cancel the singularity in
the local energy and stop the network from re-introducing one, so the local energy is bounded, the
quadrature assumptions of their generalization section actually hold, and the SR forces have finite
variance. This is the concrete sense in which the cusp engineering is conditioning, not hygiene.''
\end{say}

\vspace{2em}\hrule
\begin{center}\Large\textbf{Part III --- Synthesis: our results through both lenses}\end{center}
\hrule\vspace{1em}

% ============================================================
\begin{slide}{Slide 16 --- Best results and why CTNN beats a plain FFNN backflow (100\,s)}
\textbf{On the slide:}
\begin{itemize}[leftmargin=*]
\item VMC$+$SR energies vs DMC: $N{=}2$ within DMC error bars; $N\in\{6,12\}$ at
$10^{-2}$--$10^{-1}\%$; PINN$+$CTNN best column reaches e.g.\ $-0.0000\%$ at $(6,0.5)$,
$+0.0018\%$ at $(12,0.5)$, and extends to $N{=}20$.
\item \textbf{CTNN cell ablation} (zero message-passing, keep pairwise features), $N{=}6$:
$\Delta E=+22.4\%$ ($\omega{=}1$), $+30.5\%$ ($\omega{=}0.1$), $+11.8\%$ ($\omega{=}0.001$).
\item The entire increase is \textbf{kinetic}; potential unchanged $\Rightarrow$ the cell improves
\emph{nodal geometry}, not amplitudes.
\item Three-body sensitivity ratio: $2.03,\,2.93,\,1.05$ --- collapses in the Wigner limit.
\end{itemize}
\end{slide}

\begin{say}
``Now the results, read through both papers. First the energies: with VMC and SR we match diffusion
Monte Carlo within its error bars for two electrons and stay at hundredths-to-tenths of a percent for
six and twelve, with the CTNN variant giving the best column and extending to twenty electrons where
we have no other neural baseline.

The sharpest single experiment is the CTNN ablation. We zero the message-passing weights but keep
every pairwise feature, so the only thing removed is the graph-neural-network piece --- the
accumulated node state that encodes beyond-pairwise correlation. The energy degrades by 22 percent at
tight confinement, 30 percent at intermediate, and only 12 percent in the deep Wigner limit. Two
facts make this a clean result. The entire degradation is in the kinetic energy --- the potential is
statistically unchanged at the same sample positions --- which means the message passing is improving
the geometry of the nodal surface, not the amplitudes between nodes. That is what a graph network's
extra expressivity is supposed to buy, and we see it in the kinetic energy specifically. And the
collapse to 12 percent in the Wigner limit is physically sensible: when the electrons crystallize
into a classical lattice, the correlation is captured by pair distances plus the lattice constraint,
and the beyond-pairwise capacity is no longer needed. So the GNN advantage from Paper A's architecture
section is present exactly where the physics is genuinely many-body, and switches off where it
reduces to pairwise --- and we can read it off a single ablation table.''
\end{say}

% ============================================================
\begin{slide}{Slide 17 --- The latent manifold: what the learned dimensions are, and how Wigner changes them (100\,s)}
\textbf{On the slide:}
\begin{itemize}[leftmargin=*]
\item \textbf{Correlator:} $r_{\rm eff}\lesssim 3$ everywhere; head reconstructed from PC1 alone
($|\rho|>0.94$, $N\ge6$). A low-dimensional order-parameter model, not a generic feature extractor.
\item \textbf{What the axes encode:} linear probes show PC1--2 predict mean radius and radial variance
($R^2\approx1$ for $N{=}2$); ``how big is the dot, how much do shells fluctuate.''
\item \textbf{$\phi\leftrightarrow\psi$ transition:} $\omega\ge0.1$ pair branch $\psi$ leads PC1
(60--73\%); $\omega\le0.01$ leading axis shifts to orbital $\phi$ / global $\mathbf g$ ($\psi<15\%$).
\item \textbf{Backflow:} high-rank ($r_{\rm eff}\approx 2N{-}2$) but \emph{local} --- near-field
configurations carry $\sim6\times$ their share of $\Delta E$. In deep Wigner it \emph{switches off}
($\Delta E\sim10^{-7}$, CKA$\approx1$).
\end{itemize}
\end{slide}

\begin{say}
``This slide answers the questions directly: what are the latent dimensions, are they different in
the Wigner regime, and what conditions well. The correlator's feature space has effective rank below
three everywhere, and for six or more electrons the scalar output is essentially perfectly
reconstructed from the first principal component alone. So the network is not behaving like a generic
high-dimensional feature extractor; it is behaving like a one-to-three parameter order-parameter
model. When we run linear probes from those principal components to physical summaries, the leading
axes predict the mean radius of the dot and the variance of the radial distribution almost perfectly.
The latent dimensions are, literally, `how big is the electron cloud' and `how much do the shells
fluctuate.'

They \emph{are} different in the Wigner regime, and this is the most interesting finding. At moderate
to weak correlation, the leading axis is dominated by the pair branch --- the network's primary
degree of freedom is tracking pair correlations. At strong correlation, the leading axis shifts to
the orbital branch and the global density statistics, with the pair branch dropping below fifteen
percent. The same fixed network reorganizes which sub-network owns its principal axis as the physics
crosses over. And the input attribution tells the same story from the input side: radial distance
dominates at tight confinement, spin dominates in the Wigner limit, because there the physics is
same-spin repulsion between lattice neighbors.

Backflow is the opposite kind of object: high-rank, nearly full-dimensional, but its energetic
leverage is local --- close pairs carry several times their share of the energy correction --- and in
the deep Wigner limit it switches off entirely, with energy corrections at the $10^{-7}$ level and
feature manifolds essentially identical with and without it. The correlator alone represents the
Wigner molecule.''
\end{say}

\begin{depth}
\textbf{Tie back to conditioning and the ``ideal information flow'' question.} The user-facing way to
say this: the ideal architecture lets information flow so the optimizer finds the low-dimensional
manifold as fast as possible --- which in Paper B's language means keeping $\kappa(A)$ small along the
trajectory. Our evidence that we achieved good flow: (i) the manifold the network settles on is
genuinely low-dimensional ($r_{\rm eff}\lesssim3$), so capacity isn't wasted; (ii) the manifold can
\emph{move} between regimes (the $\phi\leftrightarrow\psi$ switch), which only a feature-learning (not
frozen-NTK) network can do, and which our compact width permits; (iii) backflow is a near-identity
ODE flow, so it adds a high-rank \emph{local} correction without destabilizing the global manifold or
inflating local-energy variance. The redundant, mostly-safe input channels are what make this flow
well-conditioned: the network always has a bounded coordinate to express each correlation in, so
gradients never have to pass through an ill-conditioned feature to reach the right manifold.
\end{depth}

% ============================================================
\begin{slide}{Slide 18 --- Synthesis and one original observation (90\,s)}
\textbf{On the slide:}

\begin{tcolorbox}[colback=yellow!10,colframe=yellow!50!black,
title=\textbf{The two papers, unified by our example},fonttitle=\bfseries]
\textbf{Paper A:} the manifold \emph{exists} --- Barron/depth/symmetry break the curse, architecture
encodes the symmetry. \quad\textbf{Paper B:} the manifold is \emph{reachable} iff $\kappa(L^\ast L\circ
TT^\ast)$ is controlled --- and almost every trick is preconditioning.
\end{tcolorbox}

\begin{tcolorbox}[colback=green!7,colframe=green!50!black,
title=\textbf{Original observation: codimension-in-pairwise},fonttitle=\bfseries]
Let $\mathcal F_k$ be the class of $k$-round permutation-invariant message-passing nets, and
$r(\omega)=\mathrm{Var}_{\rm intra}[W^{(k)}]/\mathrm{Var}_{\rm intra}[W^{(0)}]$. Then $r\to1$ when the
optimal correlator already lies in the pairwise class $\mathcal F_0$, and $r>1$ otherwise. \emph{Data:}
$r\approx1.05$ (Wigner) vs $2$--$3$ (intermediate). A variance ratio that measures the codimension of
the physical solution inside the architectural function class.
\end{tcolorbox}
\end{slide}

\begin{say}
``To close, the synthesis. The two papers are two halves of one statement. Paper A establishes that
the manifold exists --- Barron regularity, depth, and symmetry are the rigorous reasons a neural
wavefunction can beat the curse of dimensionality, and architecture is how we tell the network where
the symmetries are. Paper B establishes that the manifold is reachable only if the problem is
conditioned --- training speed is the condition number of the Hermitian square composed with the
neural tangent kernel, and essentially every successful technique is preconditioning. Our thesis is a
worked example of both halves, and especially of the second: trap units, soft features, the analytic
cusp, the variational form, and SR are all preconditioners, and we have the negative result showing
parameter-space preconditioning alone is not enough.

And one observation that I think genuinely lives between the two papers and is in neither. Define the
function class of $k$-round message-passing networks, and take the ratio of the intra-bin variance of
our full correlator to its pairwise-ablated version. That ratio is one exactly when the optimal
solution already lies in the pairwise sub-class, and greater than one when the physics genuinely needs
the graph-network expressivity. Our data gives about 1.05 in the Wigner regime and 2 to 3 at
intermediate confinement. So this variance ratio is an empirical measurement of the codimension of the
physical solution inside the architectural function class --- a way of asking, quantitatively, `does
this system actually need a graph neural network, or would pairwise suffice?' I'd offer it as a
conjecture supported by data, with a clear path to a theorem.''
\end{say}

\vspace{1em}\hrule\vspace{1em}

% ============================================================
\section*{Backup slides}

\begin{slide}{Backup A --- Local energy and REINFORCE gradient}
\[
E(\theta)=\mathbb E_{R\sim|\Psi|^2}[E_L],\quad
E_L=-\tfrac12\big(\Delta\log|\Psi|+\|\nabla\log|\Psi|\|^2\big)+V,\quad
\nabla_\theta E=2\,\mathbb E[(E_L-b)\nabla_\theta\log|\Psi|].
\]
REINFORCE differentiates only through the score $\nabla_\theta\log|\Psi|$, never through the Laplacian
path --- this is what makes backflow trainable (the catch-22 resolution).
\end{slide}

\begin{slide}{Backup B --- SR / natural gradient as preconditioner}
$S_{ij}=\mathrm{Cov}[O_i,O_j]$, $O_i=\partial_{\theta_i}\log\Psi$; update $S\,\delta\theta=-g$,
$g_i=\mathrm{Cov}[E_L,O_i]$. $S$ is the quantum Fisher matrix $=$ NTK of $\log\Psi$ $=$ the
preconditioner $P P^\top$ in Paper B's \S7.4. Fails when ESS collapses (rank-deficient $S$); Adam is
the robust fallback there.
\end{slide}

\begin{slide}{Backup C --- Barron theorem, precise statement}
$\varrho$ ReLU/sigmoidal. For $g$ with $C_g=\int\|\xi\|_2|\hat g(\xi)|d\xi<\infty$ and prob.\ measure
$\mu$ on $B_1(0)$:
$\inf_\theta\|\Phi_{(d,n,1)}(\cdot,\theta)-g\|_{L^2(\mu)}\le c\,C_g/\sqrt n$.
Dimension-free exponent; dimension hidden in $C_g$. Proof $=$ Monte Carlo over $n$ random neurons.
\end{slide}

\begin{slide}{Backup D --- Paper B condition-number theorems}
$\#(\varepsilon)=O(\kappa(A)\ln\tfrac1\varepsilon)$ (Thm 7.9); $\kappa(A)=\kappa(A\circ TT^\ast)$
with $A=L^\ast L$ (Thm 7.10). Poisson example: Fourier features, $A=D+\lambda dd^\top$, $D_{kk}=k^4$;
preconditioner $P_{kk}=1/k^2$ collapses $\kappa$. Supervised learning $=$ ($L=\mathrm{Id}$, $A=$ NTK).
\end{slide}

\begin{slide}{Backup E --- Catch-22 negative-result table (collocation backflow, $N{=}6,\omega{=}0.5$)}
Best collocation $0.32\%$; K-FAC $0.34\%$; SIR/Gumbel $0.44$--$5.0\%$; hybrid/det-filter
$2.6$--$24.8\%$; soft-Coulomb $1.41\%$. \textbf{Same ansatz, VMC$+$SR: $<0.01\%$.} No parameter-space
intervention beats the floor; only changing the form does.
\end{slide}

\vspace{1em}\hrule\vspace{1em}

% ============================================================
\section*{Anticipated questions}

\begin{anticipate}
\textbf{``Summarize Paper B's main conclusion in one sentence.''}\\
Training is the binding bottleneck for PINNs, it is governed by the condition number of the Hermitian
square $L^\ast L$ of the differential operator composed with the NTK, and almost every successful
training technique is a form of preconditioning. (Quote: ``the key bottleneck $\dots$ depends on the
spectral properties of the Hermitian square of the underlying differential operator, although suitable
preconditioning might ameliorate training issues.'')
\end{anticipate}

\begin{anticipate}
\textbf{``What is a Barron space and why does it matter for a wavefunction?''}\\
A function whose Fourier transform has a finite first moment $C_g=\int\|\xi\|\,|\hat g|\,d\xi$; such
functions are approximated by two-layer nets at the dimension-free rate $C_g/\sqrt n$. It matters
because there is now a regularity theory for the static Schr\"odinger equation in spectral Barron
spaces, so ``a neural network can represent this wavefunction efficiently'' is a theorem-backed claim,
not a hope. The cusp is the part that is \emph{not} Barron-nice, which is why we hard-wire it.
\end{anticipate}

\begin{anticipate}
\textbf{``You keep saying conditioning. Concretely, what reduces the condition number in your code?''}\\
Five things, each mapping to Paper B's $L^\ast L\circ TT^\ast$: trap units shrink $k_{\max}/k_{\min}$;
soft-core features and the analytic cusp remove the high-frequency / singular content the residual
would otherwise carry; the variational form lowers the operator order (avoids $L^\ast L$); SR is the
explicit Fisher preconditioner. And the negative control: K-FAC/Adam, which only precondition $TT^\ast$,
do not break our error floor --- because the stiffness is in $L^\ast L$.
\end{anticipate}

\begin{anticipate}
\textbf{``Why is your backflow called continuous-time, and how does that relate to the article?''}\\
Paper A shows a ResNet is an Euler discretization of an ODE $\dot\varphi=h(t,\varphi)$. Our backflow
is a near-identity bounded flow (zero-init last layer, $\tanh$, learnable scale) --- an ODE started at
the identity --- applied as message passing on a graph of particles. So it is the composition of two
of the article's architecture ideas: the continuous-time/ResNet flow and the (omitted-but-flagged)
graph neural network.
\end{anticipate}

\begin{anticipate}
\textbf{``Why three networks? Isn't Deep Sets a universal approximator for symmetric functions?''}\\
Deep Sets is universal on permutation-\emph{invariant} functions, but a particle-only Deep Sets pools
away relative geometry and so cannot represent two-body correlation at all. The pair branch is the
minimum needed for correlation; the readout combines particle, pair, and global information. So three
networks is not redundancy --- it is the smallest decomposition that is both symmetric and
correlation-capable.
\end{anticipate}

\begin{anticipate}
\textbf{``What are the latent dimensions physically, and why so few?''}\\
The correlator's effective rank is below three across the whole grid. Linear probes show the leading
axes are the mean radius and the radial fluctuation of the electron cloud --- global size and shell
breathing. So few because the ground state really is a low-dimensional object once symmetry and cusp
are removed; that is the manifold the network finds, and finding it is why ML beats basis-expansion
methods here.
\end{anticipate}

\begin{anticipate}
\textbf{``Are the latent dimensions different in the Wigner regime?''}\\
Yes. The leading principal axis switches owner: pair branch dominates at $\omega\ge0.1$, orbital and
global-density branches dominate at $\omega\le0.01$, and input attribution shifts from radial distance
to spin. Backflow's energetic role also collapses to near zero. The dimensionality stays low but its
\emph{content} reorganizes across the crossover.
\end{anticipate}

\begin{anticipate}
\textbf{``Is SR just Adam with extra steps?''}\\
No. Adam is a diagonal (per-parameter) preconditioner; SR uses the full Fisher matrix, the QMC NTK.
For the ill-conditioned Rayleigh quotient near the optimum, the off-diagonal structure is exactly what
must be inverted, and Adam stalls where SR makes progress --- except at low ESS, where the Fisher
estimate goes rank-deficient and Adam becomes the robust choice. The crossover is physically
interpretable in $(N,\omega)$.
\end{anticipate}

\vspace{1em}\hrule\vspace{1em}

\section*{Final checklist}
\begin{enumerate}[leftmargin=*]
\item \textbf{Names:} Berner/Grohs/Kutyniok/Petersen (A); De~Ryck/Mishra (B); Barron (1993);
Jacot (NTK).
\item \textbf{Paper A one-liner:} four questions (generalize / depth / curse / optimize); the curse is
broken by structure (Barron, depth, symmetry); architecture encodes the structure.
\item \textbf{Paper B one-liner:} four matching error sources; training is the bottleneck;
$\#\text{steps}\sim\kappa(L^\ast L\circ\text{NTK})$; preconditioning is the cure.
\item \textbf{The thread:} manifold exists (A) $+$ manifold reachable iff well-conditioned (B); our
work is the worked example.
\item \textbf{Three numbers:} CTNN ablation $22/30/12\%$ ($\omega=1/0.1/0.001$); correlator
$r_{\rm eff}\lesssim3$; collocation floor $0.32\%$ vs VMC$+$SR $<0.01\%$.
\item \textbf{Two equations on the whiteboard:} local energy; $\kappa(A)=\kappa(L^\ast L\circ TT^\ast)$.
\item \textbf{Strongest single point:} the K-FAC negative result --- parameter-space preconditioning
is necessary but not sufficient; operator-level stiffness ($L^\ast L$) needs a change of form.
\end{enumerate}

\end{document}
